\documentclass[12pt]{article}
\usepackage[utf8]{inputenc}
\usepackage[T1]{fontenc}
\usepackage[margin=1in]{geometry}
\usepackage{setspace}
\usepackage{xcolor}
\usepackage[colorlinks=true, urlcolor=blue!60!black, linkcolor=blue!60!black]{hyperref}
\setlength{\parindent}{2em}
\setlength{\parskip}{1.25em}
\renewcommand{\baselinestretch}{1.0}
\begin{document}
\begin{center}
{\LARGE \textit{McDonald v. City of Chicago}, 561 U.S. 742 (2010)\par}
\vspace{0.8em}
{\large Case Brief\par}
\end{center}
\vspace{1em}

\section{Memory Jogger}

In a 5--4 decision, the Supreme Court held that the Fourteenth Amendment incorporates the Second Amendment right to keep and bear arms for self-defense against state and local governments, with a four-Justice plurality relying on the Due Process Clause and Justice Thomas concurring on the Privileges or Immunities Clause (\href{https://www.courtlistener.com/opinion/149702/mcdonald-v-city-of-chicago/}{CourtListener}; \href{https://supreme.justia.com/cases/federal/us/561/742/}{Justia}).

\section{Detailed Case Facts}

Two years after \textit{District of Columbia v. Heller}, 554 U.S. 570 (2008), held that the Second Amendment protects an individual right to possess a handgun in the home for self-defense against federal infringement, four Chicago residents---Otis McDonald, Adam Orlov, Colleen Lawson, and David Lawson---challenged municipal firearms ordinances in Chicago, Illinois, and the neighboring Village of Oak Park, Illinois (\href{https://supreme.justia.com/cases/federal/us/561/742/}{Justia}; \href{https://www.courtlistener.com/opinion/149702/mcdonald-v-city-of-chicago/}{CourtListener}). Chicago Municipal Code \S~8-20-040(a) (2009) prohibited possessing any firearm without a valid registration certificate, while \S~8-20-050(c) barred registration of almost all handguns, effectively banning private handgun possession within the city (\href{https://supreme.justia.com/cases/federal/us/561/742/}{Justia}). Similarly, Oak Park Municipal Code \S\S~27-2-1 (2007) and 27-1-1 (2009) made it unlawful for private citizens to possess any firearm, expressly including handguns (\href{https://supreme.justia.com/cases/federal/us/561/742/}{Justia}).

Chicago had enacted its handgun ban in 1982 to protect residents from property loss, injury, and death from firearms (\href{https://supreme.justia.com/cases/federal/us/561/742/}{Justia}). Petitioners argued that the ban instead left law-abiding residents vulnerable to violent crime. Otis McDonald, a community activist in his late seventies living in a high-crime Chicago neighborhood, had received violent threats from drug dealers because of his neighborhood policing efforts; Colleen Lawson's home had been targeted by burglars (\href{https://supreme.justia.com/cases/federal/us/561/742/}{Justia}). Although petitioners owned handguns stored outside city limits, Chicago's ordinances prevented them from keeping those handguns in their homes for self-defense (\href{https://supreme.justia.com/cases/federal/us/561/742/}{Justia}). Because \textit{Heller} involved the District of Columbia (a federal enclave), the constitutionality of Chicago's and Oak Park's bans turned on whether the Fourteenth Amendment makes the Second Amendment binding on state and local governments (\href{https://www.law.cornell.edu/constitution-conan/amendment-2/post-heller-issues-and-application-of-second-amendment-to-states}{Cornell LII}).

\section{Procedural History}

Following \textit{Heller}, the Chicago petitioners---joined by the Illinois State Rifle Association and the Second Amendment Foundation, Inc.---filed suit in the U.S. District Court for the Northern District of Illinois seeking declaratory relief that Chicago's handgun ban and related ordinances violated the Second and Fourteenth Amendments (\href{https://supreme.justia.com/cases/federal/us/561/742/}{Justia}). Two parallel actions were filed in the same district by the National Rifle Association (NRA) and local residents challenging the Oak Park and Chicago ordinances, and all three cases were assigned to the same district judge (\href{https://supreme.justia.com/cases/federal/us/561/742/}{Justia}).

The district court dismissed the challenges in \textit{NRA, Inc. v. Oak Park}, 617 F. Supp. 2d 752 (N.D. Ill. 2008), noting that the Seventh Circuit had previously upheld a municipal handgun ban in \textit{Quilici v. Morton Grove}, 695 F.2d 261 (7th Cir. 1982), and that \textit{Heller} had expressly declined to decide incorporation (\href{https://supreme.justia.com/cases/federal/us/561/742/}{Justia}). The U.S. Court of Appeals for the Seventh Circuit affirmed in \textit{NRA, Inc. v. Chicago}, 567 F.3d 856 (7th Cir. 2009) (Easterbrook, C.J.) (\href{https://supreme.justia.com/cases/federal/us/561/742/}{Justia}; \href{https://tile.loc.gov/storage-services/service/ll/usrep/usrep561/usrep561742/usrep561742.pdf}{U.S. Reports PDF}). Although the Seventh Circuit characterized the reasoning of three nineteenth-century Supreme Court precedents---\textit{United States v. Cruikshank}, 92 U.S. 542 (1876), \textit{Presser v. Illinois}, 116 U.S. 252 (1886), and \textit{Miller v. Texas}, 153 U.S. 535 (1894)---as ``defunct,'' it held that lower federal courts remained bound by directly applicable Supreme Court precedent until the Supreme Court itself overruled it (\href{https://supreme.justia.com/cases/federal/us/561/742/}{Justia}). The Supreme Court granted certiorari on September 30, 2009, heard oral argument on March 2, 2010, and issued its decision on June 28, 2010 (\href{https://www.supremecourt.gov/docket/docketfiles/html/public/08-1521.html}{SCOTUS Docket}; \href{https://www.scotusblog.com/cases/mcdonald-v-city-of-chicago/}{SCOTUSblog}).

\section{Judicial Votes}

The Court voted \textbf{5--4} to reverse and remand the Seventh Circuit's judgment (\href{https://www.supremecourt.gov/docket/docketfiles/html/public/08-1521.html}{SCOTUS Docket}; \href{https://www.courtlistener.com/opinion/149702/mcdonald-v-city-of-chicago/}{CourtListener}):

\begin{itemize}
\item \textbf{Opinion of the Court (Parts I, II-A, II-B, II-D, III-A, and III-B):} Justice Samuel A. Alito, Jr., joined by Chief Justice John G. Roberts, Jr., and Justices Antonin Scalia, Anthony M. Kennedy, and Clarence Thomas (\href{https://www.supremecourt.gov/docket/docketfiles/html/public/08-1521.html}{SCOTUS Docket}).
\item \textbf{Plurality Opinion (Parts II-C, IV, and V):} Justice Alito, joined by Chief Justice Roberts and Justices Scalia and Kennedy (incorporating the Second Amendment via the Fourteenth Amendment's Due Process Clause and declining to disturb the \textit{Slaughter-House Cases}) (\href{https://supreme.justia.com/cases/federal/us/561/742/}{Justia}).
\item \textbf{Concurrence:} Justice Scalia filed a separate concurring opinion (joining Justice Alito's opinion in full) (\href{https://www.scotusblog.com/cases/mcdonald-v-city-of-chicago/}{SCOTUSblog}).
\item \textbf{Concurrence in Part and in the Judgment:} Justice Thomas filed an opinion joining Parts I, II-A, II-B, II-D, III-A, and III-B of Justice Alito's opinion, but rejecting Due Process incorporation in favor of the Fourteenth Amendment's Privileges or Immunities Clause (\href{https://supreme.justia.com/cases/federal/us/561/742/}{Justia}).
\item \textbf{Dissent:} Justice John Paul Stevens filed a solo dissenting opinion (\href{https://www.supremecourt.gov/docket/docketfiles/html/public/08-1521.html}{SCOTUS Docket}).
\item \textbf{Dissent:} Justice Stephen G. Breyer filed a dissenting opinion, joined by Justices Ruth Bader Ginsburg and Sonia Sotomayor (\href{https://www.supremecourt.gov/docket/docketfiles/html/public/08-1521.html}{SCOTUS Docket}).
\end{itemize}

\section{Holding}

\textbf{Judgment:} Reversed and remanded (5--4) (\href{https://www.supremecourt.gov/docket/docketfiles/html/public/08-1521.html}{SCOTUS Docket}).

\textbf{Rule Announced:} The Fourteenth Amendment makes the Second Amendment right to keep and bear arms for the purpose of self-defense---as recognized in \textit{Heller}---fully applicable to state and local governments (\href{https://supreme.justia.com/cases/federal/us/561/742/}{Justia}; \href{https://www.courtlistener.com/opinion/149702/mcdonald-v-city-of-chicago/}{CourtListener}). Under the selective incorporation framework endorsed by the four-Justice plurality, a Bill of Rights guarantee is incorporated through the Fourteenth Amendment's Due Process Clause if it is ``fundamental to our scheme of ordered liberty'' (\textit{Duncan v. Louisiana}, 391 U.S. 145, 149 (1968)) or ``deeply rooted in this Nation's history and tradition'' (\textit{Washington v. Glucksberg}, 521 U.S. 702, 721 (1997)), and once incorporated it applies to the states identically to its application against the federal government (\href{https://supreme.justia.com/cases/federal/us/561/742/}{Justia}).

\textbf{Holding vs. Dicta:} The binding five-Justice holding is that the Fourteenth Amendment incorporates the Second Amendment right to keep and bear arms for self-defense against the states and that the right is fundamental to the American scheme of ordered liberty and deeply rooted in the Nation's history and tradition (\href{https://supreme.justia.com/cases/federal/us/561/742/}{Justia}). Because Justice Thomas did not join Parts II-C, IV, and V, no five-Justice majority agreed on \textit{which clause} of the Fourteenth Amendment (Due Process vs. Privileges or Immunities) accomplishes incorporation, though four Justices adhered to Due Process incorporation and reaffirmed the \textit{Slaughter-House Cases}, 83 U.S. (16 Wall.) 36 (1873) (\href{https://www.gwlr.org/wp-content/uploads/2012/08/79-3-Cohen.pdf}{Cohen, Geo. Wash. L. Rev.}). Statements in Part IV repeating \textit{Heller}'s assurances that incorporation does not cast doubt on longstanding prohibitions on firearm possession by felons and the mentally ill, sensitive-place restrictions, or commercial sale regulations, as well as the plurality's footnote survey of remaining unincorporated Bill of Rights provisions (Third Amendment, Fifth Amendment Grand Jury Clause, Seventh Amendment civil jury trial, and Eighth Amendment Excessive Fines Clause, later incorporated in \textit{Timbs v. Indiana}), are dicta (\href{https://supreme.justia.com/cases/federal/us/561/742/}{Justia}; \href{https://digitalcommons.liberty.edu/cgi/viewcontent.cgi?article=2451&context=honors}{Liberty Univ. Honors Assessment}).

\section{Analysis of Opinions}

\subsection{Opinion of the Court and Plurality Opinion (Justice Alito)}

Writing for a five-Justice majority in Parts I, II-A, II-B, II-D, III-A, and III-B, Justice Alito traced the evolution of incorporation doctrine from \textit{Barron v. Baltimore}, 32 U.S. (7 Pet.) 243 (1833), through the \textit{Slaughter-House Cases} (1873) and \textit{Cruikshank} (1876), to the modern era of selective incorporation (\href{https://supreme.justia.com/cases/federal/us/561/742/}{Justia}). Applying the \textit{Duncan}/\textit{Glucksberg} test, the majority concluded that individual self-defense is ``the central component'' of the Second Amendment right and is deeply rooted in American history and tradition---from the 1689 English Bill of Rights and Blackstone, through the Founding-era ratification debates and early state constitutions, to Reconstruction (\href{https://supreme.justia.com/cases/federal/us/561/742/}{Justia}). The Court highlighted post-Civil War Southern efforts to disarm African American freedmen, the Freedmen's Bureau Act of 1866 (which expressly protected ``the constitutional right to bear arms''), the Civil Rights Act of 1866, the 39th Congress's Fourteenth Amendment debates, and the fact that 22 of 37 states in 1868 protected the right to keep and bear arms in their state constitutions (\href{https://supreme.justia.com/cases/federal/us/561/742/}{Justia}). The majority rejected respondents' claim that Section 1 of the Fourteenth Amendment enacted merely an antidiscrimination rule (\href{https://supreme.justia.com/cases/federal/us/561/742/}{Justia}).

Writing for a four-Justice plurality in Parts II-C, IV, and V, Justice Alito declined to reconsider the \textit{Slaughter-House Cases} or revive the Privileges or Immunities Clause, noting the lack of scholarly or litigant consensus on that clause's full scope and the decades of settled reliance on the Due Process Clause (\href{https://supreme.justia.com/cases/federal/us/561/742/}{Justia}). The plurality rejected respondents' proposal to limit incorporation to rights recognized by every ``civilized'' legal system, rejected a watered-down ``two-track'' incorporation standard for states, rejected judicial interest-balancing, and repeated \textit{Heller}'s assurances that reasonable firearms regulations remain valid (\href{https://supreme.justia.com/cases/federal/us/561/742/}{Justia}).

\subsection{Concurring Opinion (Justice Scalia)}

Justice Scalia joined Justice Alito's opinion in full, explaining that despite his originalist misgivings about substantive due process, he acquiesced in selective incorporation of enumerated Bill of Rights guarantees because it was ``both long established and narrowly limited'' (\href{https://supreme.justia.com/cases/federal/us/561/742/}{Justia}). He wrote separately to attack Justice Stevens's dissenting methodology as an unprincipled ``living Constitution'' approach that empowers unelected judges to pick and choose rights based on subjective moral judgments rather than verifiable historical traditions (\href{https://supreme.justia.com/cases/federal/us/561/742/}{Justia}). Scalia argued that while historical inquiry is imperfect, it constrains judicial discretion far better than Stevens's ``omnidirectional guideposts'' and intrudes far less on the democratic process (\href{https://supreme.justia.com/cases/federal/us/561/742/}{Justia}).

\subsection{Opinion Concurring in Part and Concurring in the Judgment (Justice Thomas)}

Justice Thomas agreed that the Fourteenth Amendment makes the Second Amendment fully binding on the states, but refused to join the plurality's Due Process rationale (\href{https://supreme.justia.com/cases/federal/us/561/742/}{Justia}). Calling substantive due process a ``legal fiction'' devoid of textual or historical grounding, Thomas argued that the Fourteenth Amendment's Privileges or Immunities Clause---not the Due Process Clause---incorporates enumerated Bill of Rights protections against the states (\href{https://supreme.justia.com/cases/federal/us/561/742/}{Justia}). Surveying English charters, antebellum treaties, the Missouri Compromise debates, speeches by Representative John Bingham and Senator Jacob Howard, Reconstruction civil rights statutes, and the violence inflicted on disarmed freedmen, Thomas concluded that the ratifying public in 1868 understood ``privileges or immunities of citizens of the United States'' to include the individual right to keep and bear arms as a minimum baseline of national citizenship, and that \textit{Cruikshank} and \textit{Slaughter-House} should be rejected (\href{https://supreme.justia.com/cases/federal/us/561/742/}{Justia}; \href{https://www.law.cornell.edu/constitution-conan/amendment-2/post-heller-issues-and-application-of-second-amendment-to-states}{Cornell LII}).

\subsection{Dissenting Opinion (Justice Stevens)}

In his final major dissent before retiring, Justice Stevens rejected both Privileges or Immunities revival (citing 137 years of precedent and the danger of a ``clean slate'') and mechanical ``jot-for-jot'' incorporation of the Second Amendment (\href{https://biz.libretexts.org/Bookshelves/Constitutional_Law/Liberty_Equality_and_Due_Process\%3A_Cases_Controversies_and_Contexts_in_Constitutional_Law/01\%3A_Chapters/1.10\%3A_CHAPTER_EIGHT-_Incorporation_and_Fundamental_Rights}{LibreTexts / Stevens Dissent}). Framing the issue as a standalone substantive due process inquiry into the Fourteenth Amendment's ``liberty clause,'' Stevens argued that rights protected against the states need not mirror federal Bill of Rights guarantees identically (\href{https://supreme.justia.com/cases/federal/us/561/742/}{Justia}). He contended that private handgun ownership is not implicit in the concept of ordered liberty because firearms have a ``fundamentally ambivalent relationship to liberty,'' gun ownership is not critical to personal autonomy or dignity in the same way as privacy or expression, other advanced democracies restrict handguns, the Second Amendment was originally a federalism provision, and states have a strong interest in local experimentation to address gun violence (\href{https://supreme.justia.com/cases/federal/us/561/742/}{Justia}; \href{https://scholarship.law.tamu.edu/context/txwes-lr/article/1481/viewcontent/10.37419_3ATWLR.V19.I3.3.pdf}{Tex. Wesleyan L. Rev.}).

\subsection{Dissenting Opinion (Justice Breyer, Joined by Justices Ginsburg and Sotomayor)}

Justice Breyer dissented on the ground that ``nothing in the Second Amendment's text, history, or underlying rationale'' warrants treating private armed self-defense as a fundamental right (\href{https://www.khanacademy.org/humanities/constitution-101/x7a03a96a83aa80ff:the-bill-of-rights/x7a03a96a83aa80ff:second-amendment-modern-cases-and-debates/a/primary-source-mcdonald-v-city-of-chicago-2012}{Khan Academy / Breyer Dissent}; \href{https://www.law.cornell.edu/constitution-conan/amendment-2/post-heller-issues-and-application-of-second-amendment-to-states}{Cornell LII}). Challenging \textit{Heller}'s historical conclusions as contested and inconclusive, Breyer urged the Court to weigh functional and structural factors: (1) the lack of popular consensus that private handgun possession is fundamental; (2) the absence of a political-process failure disadvantaging gun owners; (3) federalism and the traditional state police power over local public safety; and (4) institutional competence, warning that incorporation would force federal judges to second-guess state legislatures on complex, life-and-death empirical questions about gun regulation (\href{https://supreme.justia.com/cases/federal/us/561/742/}{Justia}; \href{https://irlaw.umkc.edu/faculty_works/671/}{Rostron, Geo. Wash. L. Rev.}).

\section{Examples: Future Applications}

\begin{enumerate}
\item \textbf{Same-Side (Invalidated / Incorporated): Municipal Ban on Handguns in Public Housing.} A city housing authority adopts a lease provision and municipal ordinance barring all tenants from possessing any operable handgun inside their apartments. Under \textit{McDonald}, because the Second Amendment applies with full force to state and local governments, a categorical ban on keeping a handgun in the home for self-defense is unconstitutional (\href{https://supreme.justia.com/cases/federal/us/561/742/}{Justia}).
\item \textbf{Same-Side (Incorporated Against States): Incorporation of the Eighth Amendment Excessive Fines Clause.} A state seizes a defendant's \$42,000 vehicle as civil forfeiture for a minor drug offense carrying a maximum \$10,000 statutory fine. Applying \textit{McDonald}'s selective incorporation test---whether the Bill of Rights guarantee is fundamental to our scheme of ordered liberty or deeply rooted in the Nation's history and tradition---the protection against excessive fines is incorporated against the states via the Due Process Clause (as the Court later held in \textit{Timbs v. Indiana}) (\href{https://digitalcommons.liberty.edu/cgi/viewcontent.cgi?article=2451&context=honors}{Liberty Univ. Honors Assessment}).
\item \textbf{Opposite-Side (Upheld / Not Invalidated): State Prohibition on Firearm Possession by Convicted Violent Felons.} A state enacts a statute making it a felony for any person previously convicted of a violent felony to possess a handgun, even inside the home. \textit{McDonald} expressly repeated \textit{Heller}'s assurance that incorporation does not cast doubt on longstanding prohibitions on the possession of firearms by felons and the mentally ill (\href{https://supreme.justia.com/cases/federal/us/561/742/}{Justia}).
\item \textbf{Opposite-Side (Not Incorporated): State Elimination of Grand Jury Indictment in Felony Cases.} A state prosecutes a felony defendant by prosecutorial information rather than by grand jury indictment under the Fifth Amendment. \textit{McDonald}'s plurality explicitly noted that pre-selective-incorporation precedents declining to incorporate the Fifth Amendment Grand Jury Clause (\textit{Hurtado v. California}) and Seventh Amendment civil jury requirement (\textit{Minneapolis \& St. Louis R. Co. v. Bombolis}) remain undisturbed by \textit{stare decisis} (\href{https://supreme.justia.com/cases/federal/us/561/742/}{Justia}).
\item \textbf{Fence-Sitter (Close / Unresolved by \textit{McDonald}): State Law Barring Lawful Permanent Residents (Non-Citizens) from Possessing Handguns.} A state prohibits lawful permanent residents from owning handguns for home self-defense. This creates a sharp doctrinal puzzle under \textit{McDonald}'s 4--1--4 split: the Due Process Clause protects ``any person'' (including non-citizens), whereas Justice Thomas's fifth vote relied on the Privileges or Immunities Clause, which textually protects only ``citizens of the United States'' (\href{https://digitalcommons.law.umaryland.edu/cgi/viewcontent.cgi?article=3523&context=mlr}{Maryland L. Rev.}; \href{https://www.gwlr.org/wp-content/uploads/2012/08/79-3-Cohen.pdf}{Cohen, Geo. Wash. L. Rev.}). Lower courts must resolve whether the plurality's Due Process rationale controls under \textit{Marks v. United States} or whether the voting paradox leaves non-citizens unprotected.
\end{enumerate}

\section{Critique}

\textbf{Originalist and Textualist Perspectives:} Originalist scholars widely praised the result in \textit{McDonald}---confirming that the Reconstruction Congress sought to protect freedmen's right to armed self-defense---while sharply criticizing the plurality's refusal to overrule the \textit{Slaughter-House Cases} (\href{https://digitalcommons.law.uga.edu/cgi/viewcontent.cgi?article=1296&context=glr}{Georgia L. Rev.}; \href{https://scholarship.law.duke.edu/cgi/viewcontent.cgi?article=1062&context=djclpp_sidebar}{Duke J. Const. L. \& Pub. Pol'y}). As scholars across the ideological spectrum note, there is near-unanimity that \textit{Slaughter-House} misread the Privileges or Immunities Clause, and relying on ``substantive due process'' perpetuates a textual oxymoron (\href{https://digitalcommons.law.uga.edu/cgi/viewcontent.cgi?article=1296&context=glr}{Georgia L. Rev.}). Yet even among originalists, debate persists over whether reviving Privileges or Immunities would unleash unenumerated economic or natural-law rights, and whether the scope of incorporated rights should be pegged to 1791 or 1868 public understanding (\href{https://scholarship.law.wm.edu/cgi/viewcontent.cgi?article=1567&context=wmborj}{Wm. \& Mary Bill Rts. J.}; \href{https://www.supremecourt.gov/DocketPDF/24/24-1046/357868/20250501090640899_24-1046\%20Amicus\%20Brief.pdf}{SCOTUS Amicus Brief}).

\textbf{Progressive and Doctrinal Perspectives:} Progressive and doctrinal critics contend that \textit{McDonald} entrenched \textit{Heller}'s contested ``law-office history,'' stripped cities of democratic tools to combat urban gun violence, and left lower courts without a workable standard of scrutiny (\href{https://irlaw.umkc.edu/faculty_works/671/}{Rostron, Geo. Wash. L. Rev.}; \href{https://scholarship.law.tamu.edu/context/txwes-lr/article/1481/viewcontent/10.37419_3ATWLR.V19.I3.3.pdf}{Tex. Wesleyan L. Rev.}). Doctrinal scholars have also highlighted the ``voting paradox'' created by the 4--1--4 split: five Justices (Thomas plus the four dissenters) rejected Due Process incorporation of the Second Amendment, while eight Justices (Alito's four plus the four dissenters) rejected Privileges or Immunities incorporation, leaving the doctrinal foundation of the holding fractured (\href{https://www.gwlr.org/wp-content/uploads/2012/08/79-3-Cohen.pdf}{Cohen, Geo. Wash. L. Rev.}). Progressive originalists further observed that the plurality's refusal to revive Privileges or Immunities signaled skepticism toward using that clause to protect unenumerated progressive rights (\href{https://scholarship.law.wm.edu/cgi/viewcontent.cgi?article=1567&context=wmborj}{Wm. \& Mary Bill Rts. J.}).

\textbf{Underlying Competing Values:} \textit{McDonald} squarely pits uniform national protection of individual liberty---specifically armed self-defense, historically vital to vulnerable minorities---against federalism, local democratic experimentation, and the state police power to reduce firearm mortality. Methodologically, it exposes the tension between \textit{stare decisis} and originalism, and between tradition-bound constitutionalism and pragmatic, evolving substantive due process (\href{https://supreme.justia.com/cases/federal/us/561/742/}{Justia}).

\section{Key Quotations}

\begin{quote}
``Self-defense is a basic right, recognized by many legal systems from ancient times to the present day, and in \textit{Heller}, we held that individual self-defense is `the central component' of the Second Amendment right.''
\end{quote}
--- Justice Alito, Opinion of the Court, 561 U.S. at 767 (\href{https://www.courtlistener.com/opinion/149702/mcdonald-v-city-of-chicago/}{CourtListener}; \href{https://supreme.justia.com/cases/federal/us/561/742/}{Justia}).

\begin{quote}
``In \textit{Heller}, we held that the Second Amendment protects the right to possess a handgun in the home for the purpose of self-defense. Unless considerations of \textit{stare decisis} counsel otherwise, a provision of the Bill of Rights that protects a right that is fundamental from an American perspective applies equally to the Federal Government and the States. See \textit{Duncan}, 391 U.S., at 149, and n.~14. We therefore hold that the Due Process Clause of the Fourteenth Amendment incorporates the Second Amendment right recognized in \textit{Heller}.''
\end{quote}
--- Justice Alito, Plurality Opinion, 561 U.S. at 791 (\href{https://www.courtlistener.com/opinion/149702/mcdonald-v-city-of-chicago/}{CourtListener}; \href{https://supreme.justia.com/cases/federal/us/561/742/}{Justia}).

\begin{quote}
``The notion that a constitutional provision that guarantees only `process' before a person is deprived of life, liberty, or property could define the substance of those rights strains credulity for even the most casual user of words.''
\end{quote}
--- Justice Thomas, Concurring in Part and Concurring in the Judgment, 561 U.S. at 811 (\href{https://www.courtlistener.com/opinion/149702/mcdonald-v-city-of-chicago/}{CourtListener}; \href{https://supreme.justia.com/cases/federal/us/561/742/}{Justia}).

\begin{quote}
``The notion that the absence of a coherent theory of the Due Process Clause will somehow curtail judicial caprice is at war with reason. Indeterminacy means opportunity for courts to impose whatever rule they like; it is the problem, not the solution.''
\end{quote}
--- Justice Scalia, Concurring, 561 U.S. at 795 (\href{https://www.courtlistener.com/opinion/149702/mcdonald-v-city-of-chicago/}{CourtListener}; \href{https://supreme.justia.com/cases/federal/us/561/742/}{Justia}).

\begin{quote}
``In sum, the Framers did not write the Second Amendment in order to protect a private right of armed self-defense.''
\end{quote}
--- Justice Breyer, Dissenting (joined by Justices Ginsburg and Sotomayor), 561 U.S. at 915 (\href{https://www.courtlistener.com/opinion/149702/mcdonald-v-city-of-chicago/}{CourtListener}; \href{https://www.khanacademy.org/humanities/constitution-101/x7a03a96a83aa80ff:the-bill-of-rights/x7a03a96a83aa80ff:second-amendment-modern-cases-and-debates/a/primary-source-mcdonald-v-city-of-chicago-2012}{Khan Academy}).

\end{document}
